\section{実軸と単位円}\label{節：実軸と単位円}
\subsection{実数の数列を単位円へ}

不動点の二解が非実数なら,\,実係数なので$\alpha,\overline\alpha$という共役の組になる.
実数$a_n$について
\[
 b_n=\frac{a_n-\alpha}{a_n-\overline\alpha}
\]
と置く.
分子と分母は共役だから,$|b_n|=1$である.
また
\[
 \kappa=\frac{p-r\alpha}{p-r\overline\alpha}
\]
も分子と分母が共役なので$|\kappa|=1$となる.
従って,$b_{n+1}=\kappa b_n$は\textbf{単位円周上の一定角の回転}である.
実軸上で複雑に見えた動きが,\,円周上では単純になる.

\noindent\textbf{【例題】}\par\nobreak
\[
 a_1=\frac12,\qquad a_{n+1}=\frac{a_n+1}{1-a_n}
\]
を複素平面で解釈し,\,周期を求めよ.

\noindent\textbf{【方針】}\par\nobreak
不動点は$i,-i$です. 差の比を取り,\,一回の回転角を調べます.

\noindent\textbf{【解答】}\par\nobreak
最小周期は4であり,
\[
 \frac12,\ 3,\ -2,\ -\frac13,\ \frac12,\ldots
\]
と繰り返す.

\noindent\textbf{【解法】}\par\nobreak
$b_n=\frac{a_n-i}{a_n+i}$とすると
\[
 b_{n+1}=\frac{1+i}{1-i}b_n=ib_n.
\]
毎回反時計回りに$\frac{\pi}{2}$回転する.
$b_n$は零でなく,4回より少ない正の回数では元へ戻らない.
復元は$a_n=i\frac{b_n+1}{1-b_n}$である.
実際に元の式で四項を計算すると上の並びを得る.
どの項も1ではないため,\,すべての更新が定義できる.

ここで,\,回転が定義できることと,\,元の実数列が全項で定義できることは区別する.
例えば同じ式で$a_1=0$とすると$a_2=1$となり,$a_3$は定義できない.
円周上では回転を続けられるが,\,復元式の分母が零になる点に到達している.

\subsection{実軸と円を結ぶ座標}\label{節：実軸と円を結ぶ座標}

回転を三角関数と結び付けるため,\,別の一定倍をした座標
\[
 w=\Phi(a)=\frac{1+ia}{1-ia}
\]
を使おう.
実数$a$なら$|w|=1$であり,\,この$\Phi(a)$は先ほどの$b$に対して$w=-b$である.
例えば$a=0$は$w=1$,$a=1$は$w=i$へ移る.

$a=\tan\theta$と書ける範囲では,
\[
 \begin{aligned}
 \Phi(\tan\theta)
 &=\frac{\cos\theta+i\sin\theta}{\cos\theta-i\sin\theta}\\
 &=\cos2\theta+i\sin2\theta.
 \end{aligned}
\]
実数$a$を傾きとして記録した角度が,\,円周上ではその2倍の角度になる.
\NewFigCayley

$w=u+iv$が単位円上で$w\ne-1$なら,\,逆に
\[
 a=\frac{v}{1+u}
\]
で戻れる.
点$(-1,0)$から$(u,v)$への直線の傾きが$a$である.
従って$\Phi$は実軸と,\,一点$-1$を除いた単位円を一対一に対応させる.
除いた一点は,\,傾きが無限に大きくなる方向に当たる.
必要に応じて実軸に無限遠点$\infty$を一つ加えると円全体と対応するが,\,実数の割り算では依然として分母の条件が必要である.

実数$t$を固定し,
\[
 f(a)=\frac{a+t}{1-ta}
\]
とすると,\,直接の計算から
\[
 \Phi(f(a))=\Phi(t)\Phi(a)
\]
となる. $t=\tan\phi$なら$\Phi(t)=e^{2i\phi}$であり,
元の式は$a=\tan\theta$の角度を$\theta+\phi$へ進める式である.
正接の加法公式
\[
 \tan(\theta+\phi)
 =\frac{\tan\theta+\tan\phi}{1-\tan\theta\tan\phi}
\]
が,\,円周上の掛け算として現れた.
角度は$\pi$の整数倍だけ違っても同じ傾きを表す.
